\section{Institutional weighting beyond sample count}
\begin{figure}[!htbp]
\centering
\resizebox{\columnwidth}{!}{\begin{tikzpicture}[
  node distance=0.5cm and 0.6cm,
  box/.style={draw, rounded corners=2pt, align=center, font=\scriptsize, minimum width=2.8cm, minimum height=0.6cm}
]
\node[box, fill=orange!8] (upd) {Client update $\Delta_i$ + metadata};
\node[box, fill=orange!12, right=of upd] (signals) {Signals\\integrity gate $\cdot$ uncertainty\\useful-uniqueness\\PathoAlign risk $\cdot$ reason codes};
\node[box, fill=orange!14, right=of signals] (weight) {Weight computation\\scoring $\cdot$ bounded caps\\mirror descent $\cdot$ conservative};
\node[box, fill=green!8, right=of weight] (agg) {Weighted aggregation};
\node[box, below=of weight] (guard) {Constraints (spec/partial)\\train/val/monitor separation\\subgroup constraints $\cdot$ $\Delta$-limits};
\draw[->] (upd) -- (signals);
\draw[->] (signals) -- (weight);
\draw[->] (weight) -- (agg);
\draw[->] (guard) -- (weight);
\end{tikzpicture}}
\caption{FAIR-WEIGHTS-H signal and constraint flow. Implemented: integrity gate,
uncertainty penalty, useful-uniqueness, bounded weights, PathoAlign audit
bridge, conservative mode. Specification-only or partial: training/validation/
monitoring weight separation, difficulty adjustment, Shapley-style contribution,
subgroup constraints, weight-change limits. No established performance or
fairness superiority (see Section~7).}
\label{fig:fair}
\end{figure}


FAIR-WEIGHTS-H is the institutional-weighting component of PathologyFL. Its purpose is to study aggregation rules in which influence can depend on more than local sample volume. The implemented engine combines adjusted validation quality, useful uniqueness, uncertainty penalties, integrity gates, bounded weights, entropy diagnostics, and conservative fallbacks.

For institution $i$, let $q_i$ denote adjusted validation quality, $u_i$ a useful-uniqueness signal, $p_i$ an uncertainty penalty, and $g_i\in\{0,1\}$ an integrity gate. The implemented score takes the form
\[
s_i = \begin{cases}
-\infty, & g_i=0,\\
\alpha q_i+\beta u_i-\lambda p_i, & g_i=1,
\end{cases}
\]
with nonnegative coefficients $\alpha,\beta,\lambda$. Scores are converted to normalized weights with either a softmax or a bounded mirror-descent update. Lower and upper weight caps prevent a single source from collapsing the aggregate, while normalized entropy and effective-client-count summaries quantify concentration of influence.

The wider design also specifies training/validation/monitoring weight separation, difficulty adjustment, subgroup constraints, contribution estimators, and explicit weight-change limits. Some of those extensions remain specification-level rather than fully implemented. Accordingly, the empirical claims in this paper come from the concrete dominance-aware PANDA experiments and the CAMELYON17 source-weighting studies, not from the method name itself or from a general fairness claim.

The scientific distinction is simple: \emph{sample count is one signal for assigning influence, not a proof that the corresponding update is the most useful for the declared objective}. The following PANDA and CAMELYON17 studies test that statement under controlled and natural center-shift settings.
